\section{Where the Consequences Return}

A robot can avoid a collision, detect damage, and request maintenance. A model can revise an answer after criticism. These are feedback capabilities. Accountability asks a further question: whose continuing position is changed by what happened?

\subsection{Body, Embodiment, and Answerability}

A body supplies a material boundary. Embodiment connects perception and action through that boundary. Answerability connects an action to continuing obligations, exposure, and response.

The distinctions matter because a machine's physical trace can be mistaken for a complete responsibility structure. A scratched casing is evidence of contact. A log is evidence that a record was made. Neither determines who compensates an injured person, investigates the incident, or carries the duty to prevent recurrence.

Self-repair can make the machine more capable \cite{terryn2017self}. It does not settle those questions. A tire bears friction; a deployer may bear the repair bill; a pedestrian may bear an injury. Several boundaries receive different consequences from one event.

The book uses \emph{consequence closure} for the degree to which action remains connected to a persistent, answerable boundary. This is a framework for practical agency and accountability. Conscious experience, moral standing, and legal personality require additional arguments; they cannot be measured by a repair log.

For a defined, finite, nonempty set of consequence items $C(a)$, one can illustrate return coverage by
\begin{equation}
 \kappa(B,a)=
 \frac{\left|\{c\in C(a):\rho_a(c)=B\}\right|}
      {|C(a)|}.
\end{equation}
$B$ is the boundary being inspected and $\rho_a$ assigns an item to its ultimate bearer in this simplified account. The result lies between zero and one.

The counting scheme matters. Dividing one injury into ten entries changes an unweighted ratio. Real consequences can be shared, uncertain, continuous, or incommensurable. A weighted or distributed assignment may be more useful in a particular investigation. The expression is an illustration of attribution, rather than a numerical test for personhood.

Complete return would also be a poor universal social goal. A driver cannot personally absorb every consequence of a collision; insurance pools some losses, emergency services provide support, and public rules distribute obligations. The aim is a justified and effective assignment, rather than forcing every effect back onto one individual.

For a living person, injury, memory, dependence, relations, and finite time connect action to an unusually persistent history. A replacement body is not an ordinary maintenance option. This gives human exposure a characteristic weight. It does not make people equally powerful, equally responsible, or equally able to escape a cost.

Machines can possess persistent identities and histories too, if an architecture maintains them. Replaceable components do not logically preclude continuity: corporations replace staff and hardware while retaining obligations. The practical issue is which boundary preserves the relevant identity and what binds it across replacement.

\subsection{The Corporation as an Engineered Boundary}

A corporation has no biological pain, yet it can own assets, undertake obligations, and face legal consequences. Law constructs a boundary through records, governance, representation, assets, and recognized continuity \cite{dewey1926historic}.

The analogy is powerful because it shows that answerability is partly an institutional achievement. Four features make the structure operational:
\begin{enumerate}
 \item An identifiable pool of resources and obligations.
 \item An identity that persists across relevant changes.
 \item Consequences that alter the boundary's future possibilities.
 \item A traceable arrangement of authority and duty.
\end{enumerate}

Actual corporate law varies, and actual corporations can externalize harm or leave victims uncompensated. Limited liability and insolvency can expose gaps. The structure supplies instruments for responsibility; its existence does not guarantee that the instruments distribute losses fairly.

An ordinary deployed model is usually a component inside such arrangements. Fining a provider affects the provider. Retraining changes the model's behavior. Withdrawing a service affects users and operators. Calling all of these “punishing the AI” conceals which boundary actually receives the consequence.

Versions need not destroy accountability. A record can preserve which version acted and a legal entity can retain obligations after changing its system. Conversely, keeping the same model name does not guarantee continuity if records and duties disappear.

The engineering and institutional jobs are complementary. Stable identifiers, versioned records, reproducible configurations, permissions, and operational logs make events investigable. Contracts, regulation, courts, insurance, and governance decide obligations and remedies.

\subsection{Distributed Action Needs Distributed Responsibility}

An output can pass through a foundation-model provider, adapter, application developer, platform, integrator, deploying organization, and end user. The contribution of each may depend on the others. A failure can arise from a combination no single component displays alone.

Many-hands problems predate AI \cite{nissenbaum1996accountability,floridi2013distributed}. Machine systems expand their speed and reach while making some causal contributions harder to inspect.

A schematic retention model writes
\begin{equation}
 R_n=R_0\prod_{i=1}^{n}\alpha_i,
 \qquad 0\leq\alpha_i\leq1,
\end{equation}
where $\alpha_i$ represents the share of a defined traceability measure retained at each handoff. It describes possible erosion when handoffs lose records or obligations. A well-designed handoff can preserve or improve the evidence, so attenuation must be established rather than assumed.

The policy question is where to attach duties. One approach identifies a party facing the affected user. Another follows control over the relevant risk. Another allocates duties across the chain and adjusts access to evidence. Their merits depend on the domain and legal setting.

Three separations are particularly dangerous: the party choosing the action differs from the party benefiting, the party holding the evidence differs from the party seeking redress, and the party assigned nominal responsibility lacks the power to change the process.

Accountability requires reconnecting these positions. A front-line employee should not become the default absorber for a system whose assumptions, update cycle, and permissions are controlled elsewhere.

The continuity that matters belongs wherever the duties can remain enforceable. It may sit with a provider, operator, professional, public body, or several of them. The model's output is the start of the inquiry, not its final answer.

This brings the problem back to people. As machines take over more production and institutions absorb more formal responsibility, the old connection between a person's skill, output, recognized position, and self-respect can loosen. The next three sections examine different ways that connection can fail.
